% !TEX root = ../main.tex

\section{Planteamiento del problema}
\label{sec:problem}

La programaci\'on, tal como se ense\~na y se practica, excluye a quien no
encaja en su molde dominante. Esto se descompone en tres tensiones
interrelacionadas que conviene examinar por separado.

\paragraph{Accesibilidad de la interfaz.}
El teclado y la pantalla son artefactos hist\'oricamente situados: fueron
dise\~nados para usuarios con destreza manual binaural, visi\'on cl\'asica
y motricidad fina sostenida. Cualquier desviaci\'on de ese perfil
como la p\'erdida temporal o permanente de movilidad en un brazo, dislexia, baja visi\'on,
edad temprana o neurodivergencia convierte la interfaz en una barrera
activa~\cite{ko23,oleson23inclusive}. Mientras que disciplinas como
la m\'usica o la pintura han desarrollado durante siglos m\'ultiples
modos de acceso, la programaci\'on se ha aferrado a un \'unico canal el cual,
aunque brinda soluciones de accesibilidad, no es intr\'insecamente inclusivo.

\paragraph{Umbral de entrada cognitivo.}
Aprender a programar requiere dominar una sintaxis
ex\'otica, una sem\'antica abstracta y un modelo mental de ejecuci\'on
que el principiante no puede inspeccionar directamente. Estudios
sistem\'aticos del proceso de aprendizaje~\cite{bosse17gerosa,robins03learning}
documentan patrones de dificultad consistentes a lo largo de instituciones distintas. 
La consecuencia es un \emph{filtro} que opera antes de que el aprendiz pueda evaluar 
si la programaci\'on le interesa.

\paragraph{Monoton\'ia paradigm\'atica.}
La mayor parte de los lenguajes actuales descienden, con variaciones de
superficie, de un mismo linaje sint\'actico. Las propuestas que
desafiaron este precepto como lenguajes visuales, en bloques o end-user
programming mejoraron la accesibilidad textual pero, en
su gran mayor\'ia, conservaron el supuesto fundamental de que el
programa se inscribe en una superficie bidimensional plana. Llevar la
programaci\'on al espacio tridimensional, al tacto y a la manipulaci\'on
de objetos sigue siendo territorio poco explorado, pese a evidencia
robusta sobre los beneficios del aprendizaje manipulativo en otras
disciplinas~\cite{montessori12mind,marshall07tangible}.

\paragraph{Pregunta de investigaci\'on.}
A partir de las tres tensiones anteriores se formula la pregunta gu\'ia:
\emph{\textquestiondown es posible dise\~nar un lenguaje de
programaci\'on tangible, Turing-completo y est\'eticamente abierto que
funcione como puerta de entrada a la computaci\'on para audiencias
m\'as diversas, sin sacrificar el rigor formal del que
depende cualquier lenguaje convencional?} El resto del art\'iculo desarrolla
una respuesta afirmativa, teniendo en cuenta distintas perspectivas y alternativas.

\endinput
